% topic_20_21.tex — Content only. Compiled via main.tex.
% ── No \documentclass, \begin{document} or \end{document} here ──

\chapteropener{20 \& 21}{Magnetic Fields \& Alternating Currents}{%
  \item Magnetic Fields and Field Lines
  \item Force on a Current-Carrying Conductor
  \item Force on a Moving Charge
  \item Hall Effect and Velocity Selector
  \item Magnetic Fields due to Currents
  \item Electromagnetic Induction
  \item Alternating Currents and r.m.s. Values
  \item Rectification and Smoothing
}

% ===================================================================
%  TOPIC 20 — MAGNETIC FIELDS
% ===================================================================

% ===== SECTION 1: CONCEPT OF MAGNETIC FIELD =====
\section{Magnetic Fields}

\begin{definitionbox}{Magnetic Field}
A \textbf{magnetic field} is a region of space in which a moving charge, or a current-carrying conductor, experiences a force.
\begin{itemize}[itemsep=2pt]
  \item Magnetic fields are produced by \textbf{moving charges} (electric currents) or by \textbf{permanent magnets}.
  \item The field is represented by \textbf{field lines} (flux lines); the direction is the direction of the force on a north pole.
  \item Field lines never cross; closer lines indicate a stronger field.
  \item Crosses ($\times$) represent field into the page; dots ($\bullet$) represent field out of the page.
\end{itemize}
\end{definitionbox}

% ===== SECTION 2: FORCE ON A CURRENT-CARRYING CONDUCTOR =====
\section{Force on a Current-Carrying Conductor}

\begin{formulabox}{Force on a Conductor}
\begin{center}
\Large $F = BIL\sin\theta$
\end{center}
\vspace{0.3em}
\begin{tabular}{@{}ll@{}}
$F$ & = force on the conductor (N) \\
$B$ & = magnetic flux density (T) \\
$I$ & = current in the conductor (A) \\
$L$ & = length of conductor in the field (m) \\
$\theta$ & = angle between the conductor and the field direction \\
\end{tabular}
\vspace{0.4em}

Force is \textbf{maximum} when $\theta = 90°$ (conductor perpendicular to field): $F = BIL$.\\
Force is \textbf{zero} when $\theta = 0°$ (conductor parallel to field).
\end{formulabox}

\begin{definitionbox}{Magnetic Flux Density $B$}
\textbf{Magnetic flux density} is defined as the force per unit current per unit length acting on a wire placed \textbf{at right angles} to the field.
$$B = \frac{F}{IL} \qquad \text{units: T (tesla) } \equiv \text{N A}^{-1}\text{m}^{-1}$$
$B$ is a vector quantity. $1\ \text{T}$ is a strong field; Earth's field is $\approx 50\ \mu\text{T}$.
\end{definitionbox}

\begin{keybox}{Fleming's Left-Hand Rule}
For a \textbf{conventional current} in a magnetic field:
\begin{itemize}[itemsep=2pt]
  \item \textbf{First finger}: direction of magnetic field $B$
  \item \textbf{Second finger}: direction of conventional current $I$
  \item \textbf{Thumb}: direction of force (motion) $F$
\end{itemize}
Remember: the rule gives the force on \emph{positive} charges / conventional current. For electrons, reverse the direction.
\end{keybox}

\begin{center}
\begin{tikzpicture}[scale=0.9]
  % Magnetic field region — crosses (into page)
  \foreach \x in {0,1,2,3,4}{
    \foreach \y in {0,1,2}{
      \node[darkgray] at (\x,\y) {$\times$};
    }
  }
  % Wire
  \draw[darkblue, very thick] (2,-0.6) -- (2,2.6);
  \node[darkblue, font=\small, anchor=west] at (2.15,2.3) {$I$ (upward)};
  % Force arrow
  \draw[->, pinkframe, very thick] (2,1) -- (0.2,1)
    node[anchor=east, font=\small, pinkframe] {$F$};
  % B label
  \node[darkgray, font=\small] at (4.5,1) {$B$ (into page)};
\end{tikzpicture}
\end{center}

% ===== SECTION 3: FORCE ON A MOVING CHARGE =====
\section{Force on a Moving Charge}

\begin{formulabox}{Force on a Moving Charge}
\begin{center}
\Large $F = BQv\sin\theta$
\end{center}
\vspace{0.3em}
\begin{tabular}{@{}ll@{}}
$F$ & = magnetic force (N) \\
$B$ & = magnetic flux density (T) \\
$Q$ & = charge (C) \\
$v$ & = speed of the charge (m s$^{-1}$) \\
$\theta$ & = angle between velocity and field \\
\end{tabular}
\vspace{0.4em}

For $\theta = 90°$: $F = BQv$, directed perpendicular to both $v$ and $B$.
\end{formulabox}

\begin{keybox}{Circular Motion in a Magnetic Field}
When a charged particle moves \textbf{perpendicular} to a uniform magnetic field, the magnetic force is always perpendicular to the velocity, so no work is done and the \textbf{speed is constant}. The particle moves in a \textbf{circle}:
$$BQv = \frac{mv^2}{r} \quad \Rightarrow \quad r = \frac{mv}{BQ}$$
A larger momentum or smaller $B$ gives a larger radius.
\end{keybox}

\subsection{The Hall Effect}

\begin{definitionbox}{Hall Voltage}
When a current-carrying conductor is placed in a magnetic field perpendicular to the current, charge carriers experience a sideways force. They accumulate on one face until the electric force balances the magnetic force, producing the \textbf{Hall voltage}:
$$V_H = \frac{BI}{ntq}$$
\begin{tabular}{@{}ll@{}}
$V_H$ & = Hall voltage (V) \\
$B$ & = magnetic flux density (T) \\
$I$ & = current through the conductor (A) \\
$n$ & = number density of charge carriers (m$^{-3}$) \\
$t$ & = thickness of the conductor in the direction of $B$ (m) \\
$q$ & = charge on each carrier (C) \\
\end{tabular}
\vspace{0.3em}

A \textbf{Hall probe} uses this effect to measure $B$: since $V_H \propto B$ at constant $I$.
\end{definitionbox}

\subsection{Velocity Selector}

\begin{keybox}{Velocity Selector}
A velocity selector uses \textbf{crossed} electric and magnetic fields so that only particles with a specific speed pass through undeflected:
$$qE = BQv \quad \Rightarrow \quad v = \frac{E}{B}$$
\begin{itemize}[itemsep=2pt]
  \item Electric force $qE$ and magnetic force $BQv$ act in \textbf{opposite} directions.
  \item Only particles where these forces balance travel in a straight line.
  \item Faster particles are deflected one way; slower particles the other.
\end{itemize}
\end{keybox}

% ===== SECTION 4: FIELDS DUE TO CURRENTS =====
\section{Magnetic Fields due to Currents}

\begin{keybox}{Field Patterns}
\begin{itemize}[itemsep=3pt]
  \item \textbf{Long straight wire}: concentric circles centred on the wire. Direction given by the \textbf{right-hand grip rule} — thumb points in direction of current, fingers curl in direction of field.
  \item \textbf{Flat circular coil}: field lines pass through the centre of the coil. Field at the centre is approximately uniform over a small region.
  \item \textbf{Long solenoid}: nearly uniform field inside, similar to a bar magnet outside. A \textbf{ferrous (iron) core} greatly increases the field strength.
\end{itemize}
\end{keybox}

\begin{center}
\begin{tikzpicture}[scale=0.85]
  % Long straight wire field
  \begin{scope}[xshift=0cm]
    \draw[darkblue, very thick] (0,-2) -- (0,2);
    \foreach \r in {0.5, 1.0, 1.5}{
      \draw[midblue, thin] (0,0) circle (\r);
    }
    % Direction arrows on circles
    \draw[->, midblue, thin] (0,0.5) -- (-0.05, 0.5);
    \draw[->, midblue, thin] (0,1.0) -- (-0.05, 1.0);
    \node[darkblue, font=\small, anchor=north] at (0,-2.3) {Straight wire};
    \node[darkblue, font=\scriptsize] at (0.3, 1.6) {$B$};
  \end{scope}
  % Solenoid field
  \begin{scope}[xshift=5.5cm]
    % Body
    \draw[darkblue, thick] (-2,-0.7) rectangle (2,0.7);
    % Internal field lines (uniform)
    \foreach \y in {-0.3, 0.0, 0.3}{
      \draw[->, midblue, thin] (-1.8,\y) -- (1.8,\y);
    }
    % External field (simplified)
    \draw[midblue, thin] (2,0) .. controls (3,1.5) and (-3,1.5) .. (-2,0);
    \draw[midblue, thin] (2,0) .. controls (3,-1.5) and (-3,-1.5) .. (-2,0);
    \node[darkblue, font=\small, anchor=north] at (0,-2.3) {Solenoid};
    \node[darkblue, font=\scriptsize, anchor=west] at (2.1,0) {N};
    \node[darkblue, font=\scriptsize, anchor=east] at (-2.1,0) {S};
  \end{scope}
\end{tikzpicture}
\end{center}

\begin{keybox}{Forces Between Parallel Conductors}
\begin{itemize}[itemsep=2pt]
  \item \textbf{Same direction currents}: conductors \textbf{attract} each other.
  \item \textbf{Opposite direction currents}: conductors \textbf{repel} each other.
  \item Each conductor sits in the magnetic field produced by the other; the force is given by $F = BIL$.
\end{itemize}
\end{keybox}

% ===== SECTION 5: ELECTROMAGNETIC INDUCTION =====
\section{Electromagnetic Induction}

\begin{definitionbox}{Magnetic Flux}
\textbf{Magnetic flux} $\Phi$ is the product of the magnetic flux density and the cross-sectional area perpendicular to the field:
$$\Phi = BA\cos\theta \qquad \text{units: Wb (weber) } \equiv \text{T m}^2$$
When $B$ is perpendicular to the area: $\Phi = BA$.\\[0.3em]
\textbf{Flux linkage} $N\Phi$ is the flux through a single turn multiplied by the number of turns $N$ of the coil: units Wb-turns.
\end{definitionbox}

\begin{formulabox}{Faraday's Law and Lenz's Law}
\begin{center}
\large $\mathcal{E} = -\dfrac{\Delta(N\Phi)}{\Delta t}$
\end{center}
\vspace{0.3em}
\textbf{Faraday's Law}: the induced e.m.f. is proportional to the rate of change of flux linkage.\\[0.3em]
\textbf{Lenz's Law}: the induced e.m.f. (and hence current) acts in a direction that \textbf{opposes} the change in flux that caused it (the minus sign above).
\vspace{0.3em}
\begin{tabular}{@{}ll@{}}
$\mathcal{E}$ & = induced e.m.f. (V) \\
$N$ & = number of turns \\
$\Delta\Phi/\Delta t$ & = rate of change of flux (Wb s$^{-1}$ $\equiv$ V) \\
\end{tabular}
\end{formulabox}

\begin{keybox}{Factors Affecting the Induced E.M.F.}
\begin{itemize}[itemsep=2pt]
  \item \textbf{Rate of change} of flux: faster change $\Rightarrow$ larger e.m.f.
  \item \textbf{Number of turns} $N$: more turns $\Rightarrow$ larger e.m.f.
  \item \textbf{Strength of field} $B$: stronger field $\Rightarrow$ larger flux change.
  \item \textbf{Area} of coil: larger area $\Rightarrow$ more flux.
\end{itemize}
Lenz's law is a consequence of conservation of energy — the induced current creates a force that opposes the motion causing induction.
\end{keybox}

\begin{warningbox}{Common Mistake}
An e.m.f. is induced only when the flux is \textbf{changing}. A conductor stationary in a steady field has zero induced e.m.f., even if the field is strong.
\end{warningbox}

% ===================================================================
%  TOPIC 21 — ALTERNATING CURRENTS
% ===================================================================

% ===== SECTION 6: ALTERNATING CURRENTS =====
\section{Alternating Currents}

\begin{definitionbox}{Characteristics of an Alternating Current}
An \textbf{alternating current} (a.c.) reverses direction periodically. For a sinusoidal a.c.:
$$I = I_0\sin\omega t \qquad V = V_0\sin\omega t$$
\begin{itemize}[itemsep=2pt]
  \item $I_0$, $V_0$: \textbf{peak} (amplitude) values.
  \item $\omega = 2\pi f = 2\pi/T$: angular frequency (rad s$^{-1}$).
  \item $T$: period (s); $f$: frequency (Hz). UK mains: $f = 50\ \text{Hz}$, $T = 20\ \text{ms}$.
\end{itemize}
\end{definitionbox}

\begin{formulabox}{R.M.S. Values}
\begin{center}
\Large $I_{\text{rms}} = \dfrac{I_0}{\sqrt{2}}$ \qquad $V_{\text{rms}} = \dfrac{V_0}{\sqrt{2}}$
\end{center}
\vspace{0.3em}
The \textbf{root-mean-square} (r.m.s.) value of an alternating current is the value of direct current that would dissipate the \textbf{same power} in a resistive load.
\end{formulabox}

\begin{formulabox}{Power in a Resistive Load}
\begin{center}
\large $P_{\text{mean}} = \tfrac{1}{2}P_{\text{max}} = \tfrac{1}{2}I_0^2 R = I_{\text{rms}}^2 R = \dfrac{V_{\text{rms}}^2}{R}$
\end{center}
\vspace{0.3em}
The mean power is \textbf{half} the peak power for a sinusoidal current.
\end{formulabox}

\subsubsection*{Sinusoidal a.c. voltage}

\begin{center}
\begin{tikzpicture}[xscale=1.1, yscale=0.9]
  \draw[->, darkblue, thick] (-0.3,0) -- (6.5,0)
    node[anchor=north, font=\small] {$t$};
  \draw[->, darkblue, thick] (0,-1.6) -- (0,1.8)
    node[anchor=east, font=\small] {$V$};
  % Sinusoid
  \draw[midblue, thick, domain=0:6.28, samples=200, smooth]
    plot (\x, {1.3*sin(180*\x/3.14)});
  % V0 label
  \draw[dashed, darkgray, thin] (0,1.3) node[anchor=east, font=\small] {$V_0$} -- (0.1,1.3);
  \draw[dashed, darkgray, thin] (0,-1.3) node[anchor=east, font=\small] {$-V_0$} -- (0.1,-1.3);
  % Vrms line
  \draw[pinkframe, dashed, thick] (0, 0.919) -- (6.28, 0.919);
  \node[pinkframe, font=\small, anchor=west] at (6.35, 0.919) {$V_{\text{rms}}$};
  % T label
  \draw[dashed, darkgray, thin] (6.28,0.05) -- (6.28,-0.25);
  \node[darkgray, font=\small, anchor=north] at (3.14,-0.05) {$T$};
  \draw[<->, darkgray, thin] (0,-0.22) -- (6.28,-0.22);
\end{tikzpicture}
\end{center}

% ===== SECTION 7: RECTIFICATION AND SMOOTHING =====
\section{Rectification and Smoothing}

\begin{definitionbox}{Rectification}
\textbf{Rectification} converts alternating current into direct current (one direction only).
\begin{itemize}[itemsep=2pt]
  \item \textbf{Half-wave rectification}: a single diode allows only one half of each cycle through. The output is a series of positive pulses with gaps.
  \item \textbf{Full-wave rectification}: a bridge rectifier (four diodes) inverts the negative half-cycles, producing a continuous pulsating d.c. output.
\end{itemize}
\end{definitionbox}

\subsubsection*{Half-wave and full-wave rectification}

\begin{center}
\begin{tikzpicture}[xscale=0.85, yscale=0.85]
  % --- Half-wave ---
  \begin{scope}
    \draw[->, darkblue, thick] (-0.2,0) -- (6.5,0)
      node[anchor=north, font=\scriptsize] {$t$};
    \draw[->, darkblue, thick] (0,-0.3) -- (0,1.7)
      node[anchor=east, font=\scriptsize] {$V$};
    % Positive half-cycles only
    \foreach \k in {0,2,4}{
      \draw[midblue, thick, domain=0:3.14, samples=100, smooth]
        plot ({\x+\k}, {1.2*sin(180*\x/3.14)});
    }
    % Gaps (zero)
    \foreach \k in {1,3,5}{
      \draw[midblue, thick] (\k,0) -- (\k+1,0);
    }
    \node[darkblue, font=\small\bfseries] at (3.2,-0.9) {Half-wave};
  \end{scope}
  % --- Full-wave ---
  \begin{scope}[yshift=-4.0cm]
    \draw[->, darkblue, thick] (-0.2,0) -- (6.5,0)
      node[anchor=north, font=\scriptsize] {$t$};
    \draw[->, darkblue, thick] (0,-0.3) -- (0,1.7)
      node[anchor=east, font=\scriptsize] {$V$};
    \foreach \k in {0,1,2,3,4,5}{
      \draw[midblue, thick, domain=0:3.14, samples=100, smooth]
        plot ({\x/2+\k/2}, {1.2*sin(180*\x/3.14)});
    }
    \node[darkblue, font=\small\bfseries] at (3.2,-0.9) {Full-wave};
  \end{scope}
\end{tikzpicture}
\end{center}

\begin{definitionbox}{Bridge Rectifier}
A bridge rectifier uses \textbf{four diodes} arranged so that:
\begin{itemize}[itemsep=2pt]
  \item During the positive half-cycle: current flows through two diodes in one path.
  \item During the negative half-cycle: current flows through the other two diodes, but in the same direction through the load.
  \item The output is always positive — both half-cycles are used.
\end{itemize}
\end{definitionbox}

\begin{center}
\begin{circuitikz}[scale=1.0, european resistors]
  % Bridge rectifier
  \draw (0,2)   to[D] (2,4);   % top-left diode
  \draw (2,0)   to[D] (4,2);   % bottom-right diode
  \draw (2,4)   to[D] (4,2);   % top-right diode
  \draw (0,2)   to[D] (2,0);   % bottom-left diode (reversed)
  % AC input on left
  \draw (0,2) -- (-1,2);
  \draw[sV] (-1,2) -- (-1,0) -- (2,0);  % not clean, use lines
  % Redefine with plain lines and source symbol
\end{circuitikz}
\end{center}

% Replace the above with a cleaner TikZ bridge diagram
\begin{center}
\begin{tikzpicture}[scale=1.1,
    every node/.style={font=\small}]
  % Nodes: top=T, bottom=Bot, left=L, right=R
  \coordinate (T)   at (2, 3);
  \coordinate (Bot) at (2, 0);
  \coordinate (L)   at (0, 1.5);
  \coordinate (R)   at (4, 1.5);

  % Diodes drawn as arrows (D1..D4)
  % D1: L -> T
  \draw[->, midblue, thick] (L) -- (T) node[midway, left] {D1};
  % D2: Bot -> L  (reversed: L -> Bot means conventional current L to Bot when forward)
  \draw[->, midblue, thick] (Bot) -- (L) node[midway, left] {D2};
  % D3: T -> R
  \draw[->, midblue, thick] (T) -- (R) node[midway, right] {D3};
  % D4: R -> Bot
  \draw[->, midblue, thick] (R) -- (Bot) node[midway, right] {D4};

  % AC source on left side L-Bot
  \draw (L) -- (-1, 1.5) node[midway, above] {a.c.};
  \draw (Bot) -- (-1, 0);
  \draw[darkblue] (-1,0) -- (-1,1.5);
  \node[darkblue] at (-1,0.75) {$\sim$};

  % Load on right side T-Bot via R node
  \draw (T) -- (4, 3) -- (5.2, 3);
  \draw (Bot) -- (4, 0) -- (5.2, 0);
  \draw[darkblue, thick] (5.2,0) -- (5.2,3);
  \node[darkblue, anchor=west] at (5.3, 1.5) {$R_L$};

  % Labels
  \fill[darkblue] (T) circle (2pt);
  \fill[darkblue] (Bot) circle (2pt);
  \fill[darkblue] (L) circle (2pt);
  \fill[darkblue] (R) circle (2pt);
\end{tikzpicture}
\end{center}

\begin{keybox}{Smoothing with a Capacitor}
A capacitor connected in \textbf{parallel with the load} smooths the rectified output:
\begin{itemize}[itemsep=2pt]
  \item The capacitor \textbf{charges} during each pulse peak.
  \item It \textbf{discharges} slowly through the load between peaks, maintaining a more constant voltage.
  \item \textbf{Larger $C$} or \textbf{larger $R_L$}: longer time constant $\tau = RC$, smoother output (less ripple).
  \item \textbf{Smaller $C$} or \textbf{smaller $R_L$}: faster discharge, larger ripple voltage.
\end{itemize}
\end{keybox}

\subsubsection*{Effect of smoothing capacitor on full-wave output}

\begin{center}
\begin{tikzpicture}[xscale=0.85, yscale=0.9]
  \draw[->, darkblue, thick] (-0.2,0) -- (7.5,0)
    node[anchor=north, font=\small] {$t$};
  \draw[->, darkblue, thick] (0,-0.3) -- (0,1.8)
    node[anchor=east, font=\small] {$V$};
  % Unsmoothed full-wave (dashed)
  \foreach \k in {0,1,2,3,4,5,6}{
    \draw[darkgray, dashed, domain=0:3.14, samples=80, smooth]
      plot ({\x/2+\k/2}, {1.2*sin(180*\x/3.14)});
  }
  % Smoothed output — decays slightly between peaks
  \draw[midblue, thick, domain=0:7.2, samples=400, smooth]
    plot (\x, {1.2 - 0.15*mod(\x,0.5)/0.5 * (1-cos(360*mod(\x,0.5)/0.5))});
  \draw[pinkframe, dashed, thick] (0,1.05) -- (7.3,1.05)
    node[anchor=west, font=\small] {$V_{\text{rms}}$};
  \node[midblue, font=\small] at (5.5,1.5) {smoothed};
  \node[darkgray, font=\small] at (5.5,0.45) {unsmoothed};
\end{tikzpicture}
\end{center}

% ===== SECTION 8: FORMULA SUMMARY =====
\section{Formula Summary Sheet}

\begin{minipage}{\linewidth}
\begin{tcolorbox}[colback=lightgold, colframe=accentgold]
\renewcommand{\arraystretch}{1.8}
\begin{tabular}{@{}lll@{}}
\toprule
\textbf{Formula} & \textbf{Quantity} & \textbf{Units} \\
\midrule
$F = BIL\sin\theta$ & Force on current-carrying conductor & N \\
$F = BQv\sin\theta$ & Force on moving charge & N \\
$r = mv/(BQ)$ & Radius of circular orbit in $B$ field & m \\
$V_H = BI/(ntq)$ & Hall voltage & V \\
$v = E/B$ & Velocity selector condition & m s$^{-1}$ \\
$\Phi = BA\cos\theta$ & Magnetic flux & Wb \\
$\mathcal{E} = -\Delta(N\Phi)/\Delta t$ & Faraday's / Lenz's law & V \\
$I = I_0\sin\omega t$ & Alternating current & A \\
$I_{\text{rms}} = I_0/\sqrt{2}$ & R.M.S. current & A \\
$V_{\text{rms}} = V_0/\sqrt{2}$ & R.M.S. voltage & V \\
$P_{\text{mean}} = \tfrac{1}{2}I_0^2 R = I_{\text{rms}}^2 R$ & Mean power in resistive load & W \\
\bottomrule
\end{tabular}
\end{tcolorbox}

\vspace{0.5em}
\begin{tcolorbox}[colback=lightblue, colframe=darkblue]
\textbf{Key values:}\quad $1\ \text{T} = 1\ \text{N A}^{-1}\text{m}^{-1}$;\quad $1\ \text{Wb} = 1\ \text{T m}^2$;\quad UK mains: $230\ \text{V}$ r.m.s., $50\ \text{Hz}$\\[0.3em]
\textbf{Peak mains voltage:} $V_0 = 230\sqrt{2} \approx 325\ \text{V}$
\end{tcolorbox}
\end{minipage}

% ===== SECTION 9: WORKED EXAMPLES =====
\section{Worked Examples}

\subsection*{Example 1 — Force on a Conductor}

\textbf{Question:} A wire of length $0.15\ \text{m}$ carries a current of $3.0\ \text{A}$ at $60°$ to a uniform field of $0.25\ \text{T}$. Calculate the force on the wire.

\begin{examplebox}{Solution}
$$F = BIL\sin\theta = 0.25 \times 3.0 \times 0.15 \times \sin 60°$$
$$F = 0.25 \times 3.0 \times 0.15 \times 0.866 = \mathbf{9.7\times10^{-2}\ \text{N}}$$
\end{examplebox}

\subsection*{Example 2 — Circular Motion of a Charged Particle}

\textbf{Question:} A proton (mass $1.67\times10^{-27}\ \text{kg}$, charge $1.6\times10^{-19}\ \text{C}$) moves at $2.0\times10^6\ \text{m s}^{-1}$ perpendicular to a field of $0.15\ \text{T}$. Calculate the radius of its circular path.

\begin{examplebox}{Solution}
$$r = \frac{mv}{BQ} = \frac{1.67\times10^{-27} \times 2.0\times10^6}{0.15 \times 1.6\times10^{-19}}$$
$$r = \frac{3.34\times10^{-21}}{2.40\times10^{-20}} = \mathbf{0.14\ \text{m}}$$
\end{examplebox}

\subsection*{Example 3 — Induced E.M.F.}

\textbf{Question:} A coil of 200 turns and area $50\ \text{cm}^2$ is in a field of $0.30\ \text{T}$ perpendicular to the plane of the coil. The field drops to zero in $0.040\ \text{s}$. Calculate the induced e.m.f.

\begin{examplebox}{Solution}
$$\Delta\Phi = B\Delta A = 0.30 \times 50\times10^{-4} = 1.5\times10^{-3}\ \text{Wb}$$
$$\mathcal{E} = N\frac{\Delta\Phi}{\Delta t} = 200 \times \frac{1.5\times10^{-3}}{0.040} = \mathbf{7.5\ \text{V}}$$
\end{examplebox}

\subsection*{Example 4 — R.M.S. and Peak Values}

\textbf{Question:} The mains supply has an r.m.s. voltage of $230\ \text{V}$ at $50\ \text{Hz}$. Find (a) the peak voltage, (b) the angular frequency and (c) the mean power dissipated in a $1.2\ \text{k}\Omega$ resistor.

\begin{examplebox}{Solution}
\textbf{(a)} \quad $V_0 = V_{\text{rms}}\sqrt{2} = 230\sqrt{2} = \mathbf{325\ \text{V}}$

\vspace{0.3em}
\textbf{(b)} \quad $\omega = 2\pi f = 2\pi\times50 = \mathbf{314\ \text{rad s}^{-1}}$

\vspace{0.3em}
\textbf{(c)} \quad $P = V_{\text{rms}}^2/R = 230^2/(1.2\times10^3) = 52900/1200 = \mathbf{44\ \text{W}}$
\end{examplebox}

% ===== SECTION 10: PRACTICE QUESTIONS =====
\section{Practice Exam Questions}

\subsection*{Section A — Short Answer Questions}

\begin{tcolorbox}[colback=white, colframe=midblue, breakable]

\begin{minipage}{\linewidth}
\textbf{Q1.} Define magnetic flux density and state its SI unit.
\par\hfill\textit{[2 marks]}
\vspace{2.5cm}
\end{minipage}

\begin{minipage}{\linewidth}
\textbf{Q2.} State Faraday's Law and Lenz's Law of electromagnetic induction.
\par\hfill\textit{[3 marks]}
\vspace{3cm}
\end{minipage}

\begin{minipage}{\linewidth}
\textbf{Q3.} Explain what is meant by the r.m.s. value of an alternating current and state why it is a more useful quantity than the peak value.
\par\hfill\textit{[3 marks]}
\vspace{3cm}
\end{minipage}

\begin{minipage}{\linewidth}
\textbf{Q4.} A wire of length $8.0\ \text{cm}$ is placed perpendicular to a field of flux density $0.40\ \text{T}$ and carries a current of $2.5\ \text{A}$. Calculate the force on the wire.
\par\hfill\textit{[2 marks]}
\vspace{2.5cm}
\end{minipage}

\begin{minipage}{\linewidth}
\textbf{Q5.} An a.c. supply has peak voltage $340\ \text{V}$. Calculate (a) the r.m.s. voltage and (b) the mean power delivered to a $680\ \Omega$ resistor.
\par\hfill\textit{[3 marks]}
\vspace{3cm}
\end{minipage}

\end{tcolorbox}

\subsection*{Section B — Longer Structured Questions}

\begin{tcolorbox}[colback=white, colframe=midblue, breakable]

\begin{minipage}{\linewidth}
\textbf{Q6.} An electron (mass $9.11\times10^{-31}\ \text{kg}$, charge $1.6\times10^{-19}\ \text{C}$) enters a uniform magnetic field of $2.0\times10^{-3}\ \text{T}$ perpendicular to the field with a speed of $5.0\times10^6\ \text{m s}^{-1}$.
\end{minipage}

\begin{enumerate}[label=(\alph*)]
  \item \begin{minipage}[t]{\linewidth}
  Calculate the radius of the circular path followed by the electron.
  \par\hfill\textit{[2 marks]}
  \vspace{3cm}
  \end{minipage}

  \item \begin{minipage}[t]{\linewidth}
  Explain why the speed of the electron remains constant even though a force acts on it.
  \par\hfill\textit{[2 marks]}
  \vspace{3cm}
  \end{minipage}

  \item \begin{minipage}[t]{\linewidth}
  The magnetic field strength is doubled. State and explain the effect on the radius of the orbit.
  \par\hfill\textit{[2 marks]}
  \vspace{3cm}
  \end{minipage}
\end{enumerate}

\begin{minipage}{\linewidth}
\textbf{Q7.} A rectangular coil of 80 turns and dimensions $4.0\ \text{cm} \times 6.0\ \text{cm}$ is placed with its plane perpendicular to a uniform magnetic field of $0.50\ \text{T}$.

\begin{enumerate}[label=(\alph*)]
  \item Calculate the flux linkage through the coil.
  \par\hfill\textit{[2 marks]}
  \vspace{3cm}

  \item The coil is rotated so that its plane becomes parallel to the field in $0.030\ \text{s}$. Calculate the mean induced e.m.f.
  \par\hfill\textit{[2 marks]}
  \vspace{3cm}

  \item State and explain the direction of the induced current using Lenz's law.
  \par\hfill\textit{[2 marks]}
  \vspace{3cm}
\end{enumerate}
\end{minipage}

\end{tcolorbox}

% ===== SECTION 11: MARK SCHEME =====
\section{Mark Scheme and Answers}

\begin{markscheme}

\textbf{Q1.} Magnetic flux density is the force per unit current per unit length on a wire placed at right angles to the field [1]; unit: tesla (T) or N A$^{-1}$ m$^{-1}$ [1].

\vspace{0.5em}\textbf{Q2.} Faraday's Law: the induced e.m.f. is proportional to the rate of change of flux linkage [1]; $\mathcal{E} = -\Delta(N\Phi)/\Delta t$ [1]. Lenz's Law: the induced e.m.f. acts in a direction to oppose the change in flux that caused it [1].

\vspace{0.5em}\textbf{Q3.} The r.m.s. value of an a.c. is the equivalent steady d.c. that would dissipate the same power in a resistive load [1]; it is more useful because power calculations use r.m.s. values directly ($P = I_{\text{rms}}^2 R$) [1]; the peak value changes continuously so it cannot directly give average power [1].

\vspace{0.5em}\textbf{Q4.} $F = BIL\sin90° = 0.40 \times 2.5 \times 0.080 = \mathbf{0.080\ \text{N}}$ [2].

\vspace{0.5em}\textbf{Q5(a).} $V_{\text{rms}} = V_0/\sqrt{2} = 340/\sqrt{2} = \mathbf{240\ \text{V}}$ [1].
\textbf{Q5(b).} $P = V_{\text{rms}}^2/R = 240^2/680 = 57600/680 = \mathbf{85\ \text{W}}$ [2].

\vspace{0.5em}\textbf{Q6(a).} $r = mv/(BQ) = (9.11\times10^{-31}\times5.0\times10^6)/(2.0\times10^{-3}\times1.6\times10^{-19})$ [1] $= 4.555\times10^{-24}/3.2\times10^{-22} = \mathbf{1.4\times10^{-2}\ \text{m}}$ [1].

\vspace{0.5em}\textbf{Q6(b).} The magnetic force is always perpendicular to the velocity [1]; a perpendicular force does no work, so kinetic energy (and hence speed) is unchanged [1].

\vspace{0.5em}\textbf{Q6(c).} Radius halves [1]; since $r = mv/BQ$, doubling $B$ halves $r$ (all other quantities constant) [1].

\vspace{0.5em}\textbf{Q7(a).} $\Phi = BA = 0.50\times(0.040\times0.060) = 0.50\times2.4\times10^{-3} = 1.2\times10^{-3}\ \text{Wb}$ [1]; flux linkage $= N\Phi = 80\times1.2\times10^{-3} = \mathbf{9.6\times10^{-2}\ \text{Wb-turns}}$ [1].

\vspace{0.5em}\textbf{Q7(b).} Final flux linkage $= 0$ (plane parallel to field); $\Delta(N\Phi) = 9.6\times10^{-2}$ [1]; $\mathcal{E} = \Delta(N\Phi)/\Delta t = 9.6\times10^{-2}/0.030 = \mathbf{3.2\ \text{V}}$ [1].

\vspace{0.5em}\textbf{Q7(c).} By Lenz's law the induced current opposes the decrease in flux [1]; so the induced current flows in a direction to maintain the original flux — i.e.\ in the direction that would make the coil act as an electromagnet attracting the field source / resisting the rotation [1].

\end{markscheme}

% ===== SECTION 12: REVISION CHECKLIST =====
\newpage
\section{Revision Checklist}

Use this checklist to track your understanding. Tick each box when you are confident:

\begin{tcolorbox}[colback=white, colframe=darkblue, breakable]
\renewcommand{\arraystretch}{1.5}
\begin{tabular}{@{} c p{10.5cm} r @{}}
\toprule
 & \textbf{Learning Objective} & \textbf{Confidence (1--3)} \\
\midrule
$\square$ & Define magnetic flux density; use \nf{F = BIL\sin\theta} & \\
$\square$ & Apply Fleming's left-hand rule to find force directions & \\
$\square$ & Use \nf{F = BQv\sin\theta} for force on a moving charge & \\
$\square$ & Derive and use \nf{r = mv/(BQ)} for circular orbit radius & \\
$\square$ & Explain and use the Hall effect; use \nf{V_H = BI/(ntq)} & \\
$\square$ & Explain the velocity selector condition \nf{v = E/B} & \\
$\square$ & Sketch field patterns for straight wire, circular coil and solenoid & \\
$\square$ & Explain forces between parallel current-carrying conductors & \\
$\square$ & Define magnetic flux \nf{\Phi = BA\cos\theta} and flux linkage \nf{N\Phi} & \\
$\square$ & State and apply Faraday's and Lenz's laws & \\
$\square$ & Distinguish peak and r.m.s. values; use \nf{I_\text{rms} = I_0/\sqrt{2}} & \\
$\square$ & Calculate mean power using r.m.s. values & \\
$\square$ & Distinguish half-wave and full-wave rectification & \\
$\square$ & Explain the role of a bridge rectifier (four diodes) & \\
$\square$ & Explain smoothing by a capacitor; describe effect of $C$ and $R$ & \\
\bottomrule
\end{tabular}
\vspace{0.5em}

\textit{Key: 1 = Need more work \quad 2 = Getting there \quad 3 = Confident}
\end{tcolorbox}

\vspace{1em}
\begin{motivationbox}
\centering
\textbf{\color{darkblue}Good luck with your revision!}\\[0.2em]
{\color{darkgray}\small Magnetic fields and electromagnetic induction underpin almost all electrical power generation and transmission. Faraday's law — one equation — explains the generator, the transformer, and the induction motor. Master that and you understand the modern world.}
\end{motivationbox}
